\documentclass[10pt,a4paper]{article}
\usepackage[utf8]{inputenc}
\usepackage[a4paper,left=3cm,right=2cm]{geometry}
\usepackage{amsmath}
\usepackage[thmmarks, amsmath, thref]{ntheorem}
\usepackage{amsfonts}
\usepackage{amssymb}
\usepackage{fancyhdr}
\pagestyle{fancy}
\usepackage{anysize}
\usepackage{graphicx}
\usepackage{float}
\usepackage{hyperref}
\usepackage{multirow}
\usepackage{enumitem}
\newtheorem{theorem}{Theorem}
\theoremsymbol{\ensuremath{\blacksquare}}
\newtheorem*{solution}{Solución:}
\usepackage{listings}
\usepackage{xcolor}

% Margenes y formalidades

\textwidth 6.25in % Margen Derecho y ancho del texto
\textheight 8.5in
\oddsidemargin 0in % Margen Izquierdo
\headheight 0in % Largo del texto
\leftmargin 1in
\parindent 0em % Salto entre la indentación (Sangría)https://www.overleaf.com/4519127669sdpptkhvsbfn
\topmargin -0.5in % Margen Superior
\parskip 0.5ex % Salto entre parrafos
\baselineskip 1pt

%\lhead{\includegraphics[scale=0.2]{logo.png}} % texto izquierda de la cabecera
%\chead{}                                      % texto centro de la cabecera
%\rhead{\includegraphics[scale=0.2]{FIG/logo-mba-2016.png}} % número de página a la derecha

\renewcommand{\headrulewidth}{0.4pt} % grosor de la línea de la cabecera
\renewcommand{\footrulewidth}{0.4pt} % grosor de la línea del pie

\begin{document}



\pagebreak
\vspace*{4mm} 

\begin{center}
\begin{huge}
\textbf{\\Ejercicios Econometría 2}
\end{Large}\\
%{\large 2023-2}\\
\end{center}

\begin{enumerate}

  \item  Se cuenta con la siguiente información para 5 individuos, donde $Y$ es la variable dependiente y $X_1$, $X_2$ son variables explicativas:

\[
\begin{array}{c|ccc}
\hline
\text{Individuo} & X_1 & X_2 & Y \\
\hline
1 & 4 & 1 & 5 \\
2 & 3 & 2 & 7 \\
3 & 5 & 3 & 10 \\
4 & 7 & 4 & 13 \\
5 & 6 & 4 & 12 \\
\hline
\end{array}
\]

\begin{enumerate}
  \item \textbf{(Modelo de Regresión Lineal Simple)} Estime el modelo $Y = \beta_0 + \beta_1\,X_1 + \varepsilon$, ignorando la variable $X_2$. Calcule los coeficientes, el $R^2$, el $R^2$ ajustado y las sumas de cuadrados (SST, SSR y SSE).


  \item \textbf{Discuta la posibilidad de sesgo} al omitir $X_2$.


  \item \textbf{(Modelo de Regresión Lineal Múltiple)} Estime el modelo $Y = \beta_0 + \beta_1 X_1 + \beta_2 X_2 + \varepsilon$, calcule el nuevo $R^2$ (y las sumas de cuadrados) y compare con el modelo simple. Hint: 
  \[
(X^TX)^{-1}
= \frac{1}{95}
\begin{bmatrix}
281 & -72 & 35\\
-72 & 34 & -35\\
35 & -35 & 50
\end{bmatrix}.
\]





\item ¿De que valor es el sesgo por variable omitida en este ejercicio? ¿Como es la correlación entre $X_1$ y $X_2$

\end{enumerate}




\item En una empresa se realizó un curso de un software importante para el desarrollo del trabajo, siendo este obligatorio con selección aleatoria para algunos trabajadores, y tambien los no escogidos aleatoriamente podían inscribirse de forma voluntaria.
Sea $Y$ el salario, $D\in\{0,1\}$ el tratamiento (tomó el curso) y $R\in\{0,1\}$ la asignación aleatoria.
Se observan 8 individuos (4 tratados, 4 no tratados). Entre los tratados, 3 son aleatorios ($R=1$); entre los no tratados, 2 son aleatorios.

\begin{center}
\begin{tabular}{|c| c |c |c|}
\hline
id & $Y$ & $D$ & $R$\\
\hline
1 & 950 & 1 & 1\\
2 & 900 & 1 & 1\\
3 & 880 & 1 & 1\\
4 & 820 & 1 & 0\\
5 & 780 & 0 & 1\\
6 & 740 & 0 & 1\\
7 & 700 & 0 & 0\\
8 & 650 & 0 & 0\\
\hline
\end{tabular}
\end{center}

\begin{enumerate}
\item Calcule la diferencia ingenua $\hat\beta_{\text{ing}}=\overline{Y}_{D=1}-\overline{Y}_{D=0}$.


\item En la submuestra $R=1$, calcule $\hat\beta_{R} = \overline{Y}_{D=1,R=1}-\overline{Y}_{D=0,R=1}$.


\item Interprete $\hat\beta$ como estimador cercano al $ATE$ (y al $ATT$ con efectos homogéneos). Calcule el sesgo de selección en este caso.

\item Discuta por qué $\hat\beta_{\text{ing}}$ puede estar sesgado cuando $D$ no es aleatorio.

\end{enumerate}

\item Queremos realizar la siguiente regresión de $Y=\beta_0+\beta_1X+\epsilon$, pero solo podemos tener una versión aproximada de $Y^*=Y+c$, con $c$ una constante desconocida, y $X^*=X+u$, siendo $u$ error clásico ($cov(X,u)=0, var(u)=\sigma^2_u$), por lo que sólo podemos tener una versión imprecisa de la regresión, $Y^* = \tilde\beta_0+\tilde\beta_1X^*+v$. Desmuestra que 
$$\hat{\tilde\beta}_1=\hat{\beta_1}\frac{Var(X)}{Var(X)+\sigma^2_u}$$

\end{enumerate}

\section{Formulario}


\begin{itemize}

\item  Estimadores MCO para RLS.
{
$$\hat{\beta}_0 = \bar{y} - \hat{\beta}_1 \bar{x} $$
$$\hat{\beta}_1 = \frac{\sum_{i=1}^{n}(x_i - \bar{x})(y_i-\bar{y})}{\sum_{i=1}^{n}(x_i - \bar{x})^2}=\frac{Cov(x,y)}{Var(x)}$$}


\item  Estimadores MCO para RLM.
{
$${\hat{\beta}} = (X^TX)^{-1} X^TY $$}


    \item \[ R^2  = 1 - \frac{SSR}{SST} \] 
\item  $$SST = \sum_{i=1}^{n} (y_i - \bar{y})^2 \quad \quad SSR = \sum_{i=1}^{n}  \hat{u}_i^2 $$
    \item \[ \bar{R}^2 =  1 - \frac{(n-1)SSR}{(n-k-1)SST} \]

    \item  $$\hat{\sigma}^2 = \frac{\sum_{i=1}^{n} \hat{u}_i^2 }{n-k-1}= \frac{SSR}{n - k - 1}$$
    \item \[
\text{Var}(\hat{\beta}|X) = \hat{\sigma}^2(X^{\top}X)^{-1}
\]
    \end{itemize}



\end{document}